\input includes/\format

\pdfinfo{
   /Title (Guanzi Pu)
   /Creator (TeX)
   /Author (Vit Brunner)
   /CreationDate (20261111233500)
}

~
\title
\parindent=0em
\baselineskip=1em

\vskip 5cm
\centerline{\cjktitle 管子譜}
\bigskip
\centerline{Guanzi Pu}
\vfill\break

\input includes/\format-motto

\parindent=1em
\bigskip
\noindent{\header Preface}
\medskip
\noindent Guanzi Pu is a large classical Chinese collection from the 17th century, first created by Guo Bailing, then edited by Tao Shiyu. Its Japanese name is Kanzufu; in Korean it's known as Gwan Ja Bo. Numerous editions exist, and this is just one among the many versions.

I resisted adding this collection for twenty years: Not all problems here are life and death; many are of the ``just eat something'' variety. It will be challenging to solve without the ability to verify solutions.

Having removed the duplicates from the original 1473 positions, this edition contains 1107 problems, which is still a lot. All positions are shown as {\bf black to play}.

\rightline{\it\me}
\rightline{\it November 2026}
\vfill\break

\input includes/\format-gotc
\input generated/\format-guanzi-pu

\bye
